\section{Proofs}
\label{app: proofs}

\paragraph{Notation and idealized assumptions.}
Recall that
\[
    \VC=\{x\in\{0,1\}^N:\sum_{i=1}^N x_i=1\}
\]
is the set of one-hot tokens, and let \(D([0,T],\VC)\) be the Skorokhod space of
cadlag paths with values in \(\VC\). For a path measure \(\mathbb Q\), write
\(\mathbb Q_t\) for its marginal at time \(t\).

We assume that the teacher exactly minimizes the expected loss and therefore
induces the true reverse path measure, \(\PP^{f^*}=\PP^*\).
We also assume that the minimization over \(f\) is performed over a sufficiently
rich model class, so the exact reverse predictor is available for every clean
distribution considered. We assume that the student path does not use
transitions that are impossible under the teacher path. If
\(\PP^\theta\not\ll\PP^*\), then
\(\DC_{\mathrm{KL}}(\PP^\theta\|\PP^*)=+\infty\), so the main theorem becomes
immediate.

\paragraph{Time convention.}
For a clean distribution \(p\), let \(\PP^p\) denote the path measure associated
with the fixed forward noising process, written in reverse-time coordinates so
that the clean sample is the terminal marginal of \(\PP^p\). Its forward-time
noised marginal at noise level \(t\) is
\[
    p_t(x_t)=\int q_t(x_t\mid x_0)\,p(dx_0).
\]
We write \(\PP^\theta=\PP^{p_\theta}\) and \(\PP^*=\PP^{p^*}\). Since time
reversal is a deterministic reparameterization of paths, the KL divergence is
unchanged by choosing either coordinate convention.

\subsection{Connection with the main text}
\label{app: connection with the main text}

DLM training can be viewed as fitting the reverse-time path measure associated
with a fixed forward noising process. Following the main text, we write this
using expected losses of the form \(\LC(f,p)\), where \(p\) is a clean-sample
distribution. For SEDD, Dynkin's formula~\citep{hanson2007applied,
campbell2022continuous} gives the
following distribution-level objective, up to terms independent of the optimized
model:
\begin{equation}
\label{eq:sedd-distribution}
\begin{aligned}
    \LC_{\mathrm{SEDD}}(f,p)
    =
    \EE_{p(x_0),t,q_t}
    \sum_{y\neq x_t}
    \langle x_t,y\rangle_{Q_t}
    \bigl(
      \langle f(x_t,t),y\rangle
      -
      \langle s(x_t,t\mid x_0),y\rangle
      \log \langle f(x_t,t),y\rangle
    \bigr).
\end{aligned}
\end{equation}
This is the SEDD instance of the unified loss \(\LC(f,p)\) used in
Equation~\eqref{eq:sedd loss}. For a generic clean distribution \(p\), define
\[
    f^p=\arg\min_f \LC_{\mathrm{SEDD}}(f,p).
\]
At this minimizer,
\begin{equation}
\label{eq:posterior-score}
    \langle f^p(x_t,t),y\rangle
    =
    \EE_{p(x_0\mid x_t)}
    \langle s(x_t,t\mid x_0),y\rangle .
\end{equation}
In particular, \(f^*=f^{p^*}\) denotes the exact teacher and
\(f^\theta=f^{p_\theta}\) denotes the exact fake model for the student
distribution.

\subsection{Proof of the theorem}
\label{app: proof of the theorem}

We prove the theorem in three steps. First, we show that the path KL equals the
SEDD inverse-distillation gap. Second, we transfer the result to MDLM and the
uniform objective, with Duo covered in the hard-token limit. Third, we use the
chain rule to show that the path KL controls the clean-data KL.

\begin{lemma}[SEDD path KL equals the inverse gap]
\label{lem:sedd-path-kl-gap}
Under the population assumptions above,
\begin{equation}
\label{eq:sedd-path-kl-gap}
    \DC_{\mathrm{KL}}(\PP^\theta\|\PP^*)
    =
    \LC_{\mathrm{SEDD}}(f^*,p_\theta)
    -
    \min_f \LC_{\mathrm{SEDD}}(f,p_\theta).
\end{equation}
\end{lemma}

\begin{proof}
The relative entropy between two continuous-time jump-process path measures can
be written as an integral over the relative jump-rate divergence. In the SEDD
parameterization this gives
\begin{equation}
\label{eq:ctmc-kl}
\begin{aligned}
    \DC_{\mathrm{KL}}(\PP^\theta\|\PP^*)
    &=
    \int_0^T
    \EE_{p_t^\theta(x_t)}
    \sum_{y\neq x_t}
    \langle x_t,y\rangle_{Q_t}
    \langle f^*(x_t,t),y\rangle
    \phi
    \left(
      \frac{\langle f^\theta(x_t,t),y\rangle}
           {\langle f^*(x_t,t),y\rangle}
    \right)dt,
\end{aligned}
\end{equation}
where \(\phi(r)=r\log r-r+1\). Following the linearization trick
of~\citet{kornilov2025universal}, for the positive ratios used below,
\[
    \phi(r)=\max_l \{rl+1-\exp(l)\}.
\]
Using the parameterization
\[
    l(y,x_t,t)
    =
    \log
    \frac{\langle f(x_t,t),y\rangle}
         {\langle f^*(x_t,t),y\rangle},
\]
and taking the maximum over admissible \(f\), Equation~\eqref{eq:ctmc-kl}
becomes
\begin{equation}
\label{eq:sedd-max}
\begin{aligned}
    \DC_{\mathrm{KL}}(\PP^\theta\|\PP^*)
    =
    \max_f
    \int_0^T
    \EE_{p_t^\theta(x_t)}
    \sum_{y\neq x_t}
    \langle x_t,y\rangle_{Q_t}
    \biggl[
      &\langle f^\theta(x_t,t),y\rangle
      \log
      \frac{\langle f(x_t,t),y\rangle}
           {\langle f^*(x_t,t),y\rangle}
      \\
      &+
      \langle f^*(x_t,t)-f(x_t,t),y\rangle
    \biggr]dt.
\end{aligned}
\end{equation}
We now rewrite the expectation in Equation~\eqref{eq:sedd-max}. First,
apply the posterior identity~\eqref{eq:posterior-score} with
\(p=p_\theta\):
\[
    \langle f^\theta(x_t,t),y\rangle
    =
    \EE_{p_\theta(x_0\mid x_t)}
    \langle s(x_t,t\mid x_0),y\rangle .
\]
After substituting this identity, the only term depending on \(x_0\) is
\(\langle s(x_t,t\mid x_0),y\rangle\). The remaining term
\(\langle f^*(x_t,t)-f(x_t,t),y\rangle\) depends only on \(x_t\), \(t\), and
\(y\), so it can be placed inside the same conditional expectation:
\[
\begin{aligned}
&\EE_{p_t^\theta(x_t)}
\sum_{y\neq x_t}
\langle x_t,y\rangle_{Q_t}
\biggl[
  \langle f^\theta(x_t,t),y\rangle
  \log
  \frac{\langle f(x_t,t),y\rangle}
       {\langle f^*(x_t,t),y\rangle}
  +
  \langle f^*(x_t,t)-f(x_t,t),y\rangle
\biggr]
\\
&=
\EE_{p_t^\theta(x_t)}
\EE_{p_\theta(x_0\mid x_t)}
\sum_{y\neq x_t}
\langle x_t,y\rangle_{Q_t}
\biggl[
  \langle s(x_t,t\mid x_0),y\rangle
  \log
  \frac{\langle f(x_t,t),y\rangle}
       {\langle f^*(x_t,t),y\rangle}
  +
  \langle f^*(x_t,t)-f(x_t,t),y\rangle
\biggr].
\end{aligned}
\]
Finally, the joint law of \((x_0,x_t)\) can be written in two equivalent ways:
\[
    p_t^\theta(x_t)p_\theta(x_0\mid x_t)
    =
    p_\theta(x_0)q_t(x_t\mid x_0),
\]
so
\[
    \EE_{p_t^\theta(x_t)}
    \EE_{p_\theta(x_0\mid x_t)}[\cdot]
    =
    \EE_{p_\theta(x_0),\,x_t\sim q_t(\cdot\mid x_0)}[\cdot].
\]
The maximization over \(f\) stays outside this change of variables. Therefore
we obtain
\begin{equation}
\label{eq:sedd-population-recognition}
\begin{aligned}
    \DC_{\mathrm{KL}}(\PP^\theta\|\PP^*)
    =
    \max_f
    \EE_{p_\theta(x_0),t,q_t}
    \sum_{y\neq x_t}
    \langle x_t,y\rangle_{Q_t}
    \biggl[
      &\langle s(x_t,t\mid x_0),y\rangle
      \log
      \frac{\langle f(x_t,t),y\rangle}
           {\langle f^*(x_t,t),y\rangle}
      \\
      &+
      \langle f^*(x_t,t)-f(x_t,t),y\rangle
    \biggr].
\end{aligned}
\end{equation}
The right-hand side is exactly the population SEDD loss difference:
\[
    \max_f
    \left\{
      \LC_{\mathrm{SEDD}}(f^*,p_\theta)
      -
      \LC_{\mathrm{SEDD}}(f,p_\theta)
    \right\}.
\]
Because the first term is independent of \(f\), this equals
\[
    \LC_{\mathrm{SEDD}}(f^*,p_\theta)
    -
    \min_f \LC_{\mathrm{SEDD}}(f,p_\theta).
\]
This proves the claim. Importantly, the maximum remains outside the population
expectation throughout the argument.
\end{proof}

\begin{lemma}[Gap preservation under equivalent losses]
\label{lem:gap-preservation}
Suppose two losses satisfy
\[
    \LC_A(Tf,p)=\LC_B(f,p)+C(p),
\]
where \(C(p)\) is independent of the optimized model \(f\), and \(T\) maps the
candidate class for \(\LC_B\) onto the corresponding candidate class for
\(\LC_A\). Then their inverse gaps agree:
\[
    \LC_A(Tf^*,p)-\min_g \LC_A(g,p)
    =
    \LC_B(f^*,p)-\min_f \LC_B(f,p).
\]
\end{lemma}

\begin{proof}
The claim follows by cancellation of the model-independent constant:
\[
\begin{aligned}
    \LC_A(Tf^*,p)-\min_g\LC_A(g,p)
    &=
    \LC_B(f^*,p)+C(p)
    -
    \min_f\{\LC_B(f,p)+C(p)\}
    \\
    &=
    \LC_B(f^*,p)-\min_f\LC_B(f,p).
\end{aligned}
\]
\end{proof}

\paragraph{MDLM, UDLM, and Duo.}
For the absorbing diffusion matrix \(Q_{\mathrm{abs}}\) in
Equation~\eqref{eq: absorbing diffusion matrix}, the SEDD objective under the
clean-token parameterization is equivalent to the MDLM loss
\(\LC_{\mathrm{MDLM}}\) up to terms independent of the optimized predictor.
Lemma~\ref{lem:gap-preservation} therefore gives
\[
    \DC_{\mathrm{KL}}(\PP^\theta\|\PP^*)
    =
    \LC_{\mathrm{MDLM}}(f^*,p_\theta)
    -
    \min_f \LC_{\mathrm{MDLM}}(f,p_\theta).
\]
For the uniform diffusion matrix \(Q_{\mathrm{uni}}\) in
Equation~\eqref{eq: uniform diffusion matrix}, the same argument transfers the
SEDD gap to the hard-token UDLM objective. In distribution-level notation,
\[
    \LC_{\mathrm{UDLM}}(f,p)
    \vcentcolon=
    \EE_{p(x_0),t,x_t\sim q_t(\cdot\mid x_0)}
    \bigl[
      g_{\mathrm{UDLM}}\bigl(x_t,x_0,f(x_t,t)\bigr)
    \bigr].
\]
Thus
\[
    \DC_{\mathrm{KL}}(\PP^\theta\|\PP^*)
    =
    \LC_{\mathrm{UDLM}}(f^*,p_\theta)
    -
    \min_f \LC_{\mathrm{UDLM}}(f,p_\theta).
\]
Duo evaluates the same uniform-process objective on a Gaussian-relaxed input and
recovers the hard-token UDLM objective as \(\tau\to0^+\). Hence the theorem
applies to Duo in this hard-token limit, as stated in the main text. This
argument does not claim exact uniqueness for finite-temperature Duo.

\begin{theorem}[Unique solution]
\label{thm:unique-solution-appendix}
For SEDD, MDLM, and hard-token UDLM, and for Duo in the limit
\(\tau\to0^+\), the IDLM loss satisfies
\[
    \LC_{\mathrm{IDLM}}(\theta)
    =
    \DC_{\mathrm{KL}}(\PP^\theta\|\PP^*)
    \geq
    \DC_{\mathrm{KL}}(p_\theta\|p^*)
    \geq 0.
\]
Moreover,
\[
    \LC_{\mathrm{IDLM}}(\theta)=0
    \quad\Longleftrightarrow\quad
    p_\theta=p^*.
\]
\end{theorem}

\begin{proof}
Disintegrate the path measures with respect to their clean-sample marginal:
\[
    \PP^\theta(dx,dw)=p_\theta(dx)\PP_x^\theta(dw),
    \qquad
    \PP^*(dx,dw)=p^*(dx)\PP_x^*(dw).
\]
The chain rule for relative entropy gives
\[
    \DC_{\mathrm{KL}}(\PP^\theta\|\PP^*)
    =
    \DC_{\mathrm{KL}}(p_\theta\|p^*)
    +
    \EE_{x\sim p_\theta}
    \DC_{\mathrm{KL}}(\PP_x^\theta\|\PP_x^*).
\]
The conditional KL term is nonnegative, so
\[
    \DC_{\mathrm{KL}}(\PP^\theta\|\PP^*)
    \geq
    \DC_{\mathrm{KL}}(p_\theta\|p^*)
    \geq 0.
\]
Now consider the zero-loss condition. From the previous part of the proof,
\(\LC_{\mathrm{IDLM}}(\theta)=\DC_{\mathrm{KL}}(\PP^\theta\|\PP^*)\). Therefore,
if \(\LC_{\mathrm{IDLM}}(\theta)=0\), then
\[
    \DC_{\mathrm{KL}}(\PP^\theta\|\PP^*)=0.
\]
Substituting this into the chain-rule decomposition gives
\[
    0
    =
    \DC_{\mathrm{KL}}(p_\theta\|p^*)
    +
    \EE_{x\sim p_\theta}
    \DC_{\mathrm{KL}}(\PP_x^\theta\|\PP_x^*).
\]
Both terms on the right are nonnegative, so each term must be zero. In
particular,
\[
    \DC_{\mathrm{KL}}(p_\theta\|p^*)=0.
\]
The KL divergence between two probability distributions is zero only when the
two distributions are equal, hence \(p_\theta=p^*\).

Conversely, assume \(p_\theta=p^*\). The forward noising process is fixed: once
the clean distribution \(p\) is chosen, the whole path measure \(\PP^p\) is
obtained by applying the same noising dynamics to samples from \(p\). Therefore,
if the clean distributions are equal, the induced path measures are also equal:
\[
    \PP^\theta=\PP^{p_\theta}=\PP^{p^*}=\PP^*.
\]
Thus
\[
    \DC_{\mathrm{KL}}(\PP^\theta\|\PP^*)=0,
\]
and using again
\(\LC_{\mathrm{IDLM}}(\theta)=\DC_{\mathrm{KL}}(\PP^\theta\|\PP^*)\), we get
\(\LC_{\mathrm{IDLM}}(\theta)=0\).
\end{proof}

\paragraph{Sequence-level extension.}
\label{app:sequence-level-extension}
The proof above is for one token. For a full sequence, the same proof works in
an ideal setting where the whole sequence is treated as one state and the models
describe the full reverse process over sequences. In that case, the IDLM loss is
the KL divergence between the full sequence path measures, so it also controls
the difference between \(p_\theta(x_0^{1:L})\) and \(p^*(x_0^{1:L})\). This is
not exactly the loss used in experiments. In practice, we sum token-level losses,
which is convenient for training but does not by itself give the full sequence
path KL, because tokens may still depend on each other. Thus, the proof covers
the one-token objective and the ideal full-sequence objective, but not the
practical factorized sequence loss. We leave a formal theoretical analysis of
the factorized sequence-level IDLM objective as a promising direction for future
research.
